\chapter{Objetivo e Metodologia}

O presente estudo tem como objetivo principal identificar e propor medidas para a eficientização do uso da energia elétrica na unidade consumidora localizada em <<localidade>>, com vistas à redução dos custos operacionais e à eventual adoção de fontes alternativas de energia, quando técnica e economicamente viável.

O diagnóstico energético considera, de forma integrada, os aspectos técnicos e operacionais da unidade, abrangendo instalações físicas, equipamentos, rotinas de funcionamento, práticas de manutenção e hábitos de consumo dos usuários.

\section{Objetivos Específicos}

\begin{itemize}
    \item Identificar os principais consumidores de energia elétrica e eventuais desperdícios.
    \item Avaliar o desempenho energético da unidade com base em indicadores técnicos.
    \item Propor medidas de eficiência energética com embasamento técnico e econômico.
    \item Estimar os investimentos necessários para a implementação das medidas sugeridas.
    \item Analisar a viabilidade técnico-econômica das propostas com base em critérios de custo-benefício.
    \item Contribuir para a modernização e sustentabilidade da gestão energética da unidade.
\end{itemize}

\section{Metodologia}

A metodologia adotada neste diagnóstico energético segue práticas consagradas na área de eficiência energética, conforme diretrizes do \ac{PEE} da \ac{ANEEL}~\cite{aneel920}. O processo foi desenvolvido por meio das seguintes etapas:

\begin{enumerate}
    \item \textbf{Inspeção técnica local:} Visita à unidade consumidora para reconhecimento das instalações elétricas, identificação dos ambientes e observação das condições operacionais.
    
    \item \textbf{Entrevistas:} Coleta de informações junto a funcionários e usuários, com foco em hábitos de uso, funcionamento dos sistemas e percepção de conforto ambiental.
    
    \item \textbf{Cadastramento das cargas:} Levantamento e registro de uma amostra representativa de sistemas e equipamentos elétricos instalados na unidade.
    
    \item \textbf{Análise de dados operacionais:} Avaliação do histórico de consumo, faturas de energia elétrica, demandas contratadas, principais cargas e estrutura tarifária.
    
    \item \textbf{Proposição de medidas:} Identificação de oportunidades para redução de consumo e custos, considerando tecnologias mais eficientes, ajustes operacionais e ações institucionais.
    
    \item \textbf{Estimativa de investimentos:} Cálculo dos custos médios de mercado para implementação das medidas, incluindo tributos e custos indiretos associados.
    
    \item \textbf{Avaliação de viabilidade:} Análise técnico-econômica das propostas com base em indicadores como a Relação Custo-Benefício, tempo de retorno, Valor Presente Líquido, Taxa Interna de Retorno e economia anual estimada.
    
    \item \textbf{Elaboração do relatório técnico:} Consolidação das análises, resultados, conclusões e recomendações, apresentadas de forma clara e objetiva para subsidiar a tomada de decisão.
\end{enumerate}

A abordagem metodológica adotada visa garantir elevada confiabilidade aos dados processados, fundamentando as recomendações em medições de campo, evidências técnicas e boas práticas. As medidas propostas priorizam a viabilidade de execução com investimentos justificados técnica e economicamente. As economias previstas foram estimadas com base na média dos últimos 12 meses de consumo da unidade. 
\\\\\\
{\textcolor{red}{\textbf{<<faturasDisponiveis>>}}}




%\section{Primeira Seção do Terceiro Capítulo}

